% Standalone section fragment; requires amsmath, amssymb, amsthm, longtable.
\section{The coefficient-support lemma and an explicit recomputation of LMR}
\subsection{S-001. The coefficient-support lemma, including every copy}
Let $n,d\geq1$, and let $E$ be a complex vector space with ordered basis $e_1,\ldots,e_r$.
Use the upper-triangular highest-weight convention, so $E_{i,i+1}$
sends $e_{i+1}$ to $e_i$.  A polynomial in
$\operatorname{Sym}^d(\operatorname{Sym}^n E)$ is written in the algebraically
independent generators $c_\alpha=e_1^{\alpha_1}\cdots e_r^{\alpha_r}$,
$|\alpha|=n$.

\paragraph{Statement.}
If $P$ is any highest-weight vector of weight $\lambda\vdash nd$, every
generator $c_\alpha$ in every monomial occurring with nonzero coefficient
in $P$ satisfies
\[
             \sum_{i=j}^r\alpha_i\leq\lambda_j
                   \qquad(1\leq j\leq r).
\]
For $j=1$ the inequality is automatic whenever the highest-weight space
is nonzero. In particular the statement requested as F-016 holds for
all $j\geq2$, for the complete multiplicity space.

\paragraph{Source correction.}
The attached latest GCT\_KRON paper already contains a numbered finding
F-016 and this support argument: \texttt{GCT\_KRON\_work/GCT\_KRON.tex},
subsection ``Finding F-016: tail-sum support and full determinant
multiplicity'', equivalently
\texttt{agents/formulas/nested\_support.tex}, first proof.
The intake's statement that the numbered full proof is absent is therefore
outdated for the supplied checkpoint. The following is an independent
complete proof of the requested assertion, not a claim that the lemma was
first proved in this continuation.

\begin{proof}
Introduce the matrix-coordinate algebra
$A=\operatorname{Sym}(E\otimes\mathbb C^d)$, with generators $x_{ih}$,
$1\leq i\leq r$, $1\leq h\leq d$. Its subspace of column multidegree
$(n,\ldots,n)$ is canonically $(\operatorname{Sym}^nE)^{\otimes d}$.
Symmetrizing the columns embeds $\operatorname{Sym}^d(\operatorname{Sym}^nE)$
in this subspace, equivariantly for $GL(E)$. The image of a monomial in
the $c_\alpha$ is a nonzero scalar times the sum of its distinct column
assignments. Two distinct coefficient monomials have disjoint sets of
column assignments.

By the Cauchy decomposition
\[
 A_{nd}=\bigoplus_{\mu\vdash nd}S_\mu E\otimes S_\mu\mathbb C^d,
\]
the row highest-weight space of weight $\lambda$ is one copy of the
irreducible column module $S_\lambda\mathbb C^d$. If this module is zero,
there is no assertion to prove. Otherwise $\ell(\lambda)\leq d$.
For $1\leq h\leq\ell(\lambda)$ let $\Delta_h$ be the minor in the
first $h$ rows and first $h$ columns. The nonzero polynomial
\[
                 \prod_{h\geq1}\Delta_h^{\lambda_h-\lambda_{h+1}}
\]
is row highest-weight of weight $\lambda$. Its $GL_d$ column span is the
whole row highest-weight space by irreducibility. Under a column change,
Cauchy--Binet expresses a height-$h$ flag minor as a linear combination
of minors with the same first $h$ rows. Consequently the entire row
highest-weight space is spanned by products of flag minors whose heights
are the column lengths of $\lambda$.

Exactly $\lambda_j$ factors of each spanning product have height at
least $j$. A minor of smaller height has no variable from rows
$j,\ldots,r$. A determinant monomial in any other minor has at most one
variable in a fixed matrix column. Thus every matrix monomial in every
spanning product has degree at most $\lambda_j$ in the entries of rows
$j,\ldots,r$ of each fixed column. Linear combination preserves this
support property. Projecting onto column multidegree $(n,\ldots,n)$ also
preserves it, since that projection only deletes monomials.

Apply this bound to the symmetrization of $P$. If one coefficient
monomial of $P$ contained a forbidden generator, every column assignment
containing that generator would violate the matrix-column bound. Its
coefficient in the symmetrization is nonzero and cannot cancel with a
different coefficient monomial, by the disjoint-support observation.
This contradiction proves the bound for every monomial of $P$.
The argument applies to an arbitrary vector in the full highest-weight
space, so it does not choose or discard multiplicity copies.
\end{proof}

\paragraph{Dictionary and scope.}
This proof is a calculation with exterior powers of a universal map of
locally free sheaves and its torus weight grading. Its representation
ring shadow is the Cauchy identity; that shadow alone does not establish
coefficient support or the rank of a coherent quotient. In the supplied
Grothendieck archive transcription, cote~133, scan~028, the Cauchy tensor
power identity and the question of an actual direct summand are the
relevant exact passage. The affine group-scheme representation/comodule
and diagonalizable-group grading dictionary is SGA~3, Expos\'e~I,
Propositions~4.7.2--4.7.3. The use of irreducibility and symmetrization
above is on the characteristic-zero fiber. No positive-characteristic
semisimplicity, purity, or motivic assertion is inferred.

\subsection{S-002. Six explicit degree-twelve highest-weight functions}
Let $V=\mathbb C^7$ with basis $e_0,\ldots,e_6$, and write a cubic
$F(v)=T(v,v,v)$ with symmetric trilinear form $T$. Put
\[
 a=T_{000},\quad b=T_{001},\quad c=T_{011},\quad d=T_{111},
 \qquad S_i=(T_{ijk})_{0\leq j,k\leq6}\quad(i=0,1),
 \qquad A=(T_{00j})_{0\leq j\leq6}.
\]
Here $d$ in this subsection is a binary cubic coefficient; the outer
polynomial degree is always twelve. Define
\[
 \begin{split}
 H(X,Y)&=\det(XS_0+YS_1)
       =\sum_{k=0}^7\binom7k h_kX^{7-k}Y^k,\\
 B(X,Y)&=\det\begin{pmatrix}XS_0+YS_1&A\\A^t&0\end{pmatrix}
       =\sum_{k=0}^6\binom6k b_kX^{6-k}Y^k.
 \end{split}
\]
The lower-case $b$ without a subscript is $T_{001}$; $b_k$ denotes a
bordered-determinant coefficient. The following six explicit functions
have degree twelve in $T$:
\[
\begin{aligned}
 P_0={}&h_0(2b^5-5ab^3c+3a^2bc^2+a^2b^2d-a^3cd),\\
 P_1={}&h_0(b^5-2ab^3c+a^2bc^2)
        +h_1(-ab^4+2a^2b^2c-a^3c^2),\\
 P_2={}&h_0(2b^5-3ab^3c+a^2b^2d)
        +h_1(-4ab^4+6a^2b^2c-2a^3bd)\\
       &+h_2(2a^2b^3-3a^3bc+a^4d),\\
 P_3={}&h_0(b^5-ab^3c)+h_1(-3ab^4+3a^2b^2c)\\
       &+h_2(3a^2b^3-3a^3bc)+h_3(-a^3b^2+a^4c),\\
 P_4={}&h_0b^5-5h_1ab^4+10h_2a^2b^3-10h_3a^3b^2
        +5h_4a^4b-h_5a^5,\\
 R={}&b_0(5b^3c-3abc^2-ab^2d+a^2cd)-10b_1b^4
       +20b_2ab^3-20b_3a^2b^2+10b_4a^3b-2b_5a^4.
\end{aligned}
\]
They are highest-weight vectors of weight
\[
                         \lambda=(19,7,2,2,2,2,2).
\]
They form a basis of the entire ambient highest-weight space. In
particular $R$ supplies a copy outside the degree-twelve binary
cubic--Hessian-determinant span of $P_0,\ldots,P_4$.

\begin{proof}[Highest-weight calculation]
Let $g$ be upper triangular. Polynomial functions in
$\operatorname{Sym}^{12}(\operatorname{Sym}^3V)$ are evaluated with the
convention $(gP)(F)=P(F\circ g)$. If $g$ fixes $e_0,e_1$ pointwise,
then $S_i(F\circ g)=g^tS_i(F)g$ and $A(F\circ g)=g^tA(F)$.
Both determinant constructions therefore acquire the common factor
$\det(g)^2$. Thus the unipotent radical fixing $e_0,e_1$ fixes all six
functions.

For the remaining simple raising operator $D=E_{0,1}$, its action is
\[
 Da=0,\quad Db=a,\quad Dc=2b,\quad Dd=3c,
 \qquad Dh_k=kh_{k-1},\qquad Db_k=kb_{k-1}.
\]
For the last identity $A$ is fixed because $e_0$ is fixed; the coefficient
normalizations by $\binom7k$ and $\binom6k$ give exactly the displayed
multipliers. Substitution and the product rule give $DP_i=DR=0$.
The standalone checker verifies these six identities by exact sparse
integer arithmetic, with no evaluations or probabilistic step.

Under a diagonal $g=\operatorname{diag}(t_0,\ldots,t_6)$,
$h_k$ has character $\det(g)^2t_0^{7-k}t_1^k$ and $b_k$ has character
$\det(g)^2t_0^{10-k}t_1^k$. Each displayed summand therefore has character
$\det(g)^2t_0^{17}t_1^5$, which is the claimed partition. The $h_k$ have
coefficient degree seven and the $b_k$ degree eight, proving degree twelve.
This establishes the highest-weight assertion directly, without
coefficient-support assumptions.
\end{proof}

\begin{proof}[Ambient dimension and independence]
The ambient multiplicity is independently evaluated over $\mathbb Q$ by
\texttt{check\_lmr\_ambient.py}. Its complete mathematical prescription is
\[
 h_0[h_3]=1,\qquad
 k h_k[h_3]=\sum_{i=1}^k h_{k-i}[h_3]
          \frac{p_i^3+3p_i p_{2i}+2p_{3i}}6,
 \qquad
 a_\lambda=\langle h_{12}[h_3],s_\lambda\rangle.
\]
The code expands the recurrence in power sums using exact rational
arithmetic and computes $\langle p_\rho,s_\lambda\rangle=\chi^\lambda(\rho)$
by Murnaghan--Nakayama beta-number moves. There are 3560 nonzero
power-sum types; the result is exactly six. Independently of the final
value, the cycle-index mass is checked to equal one and the identity
character is checked against the hook-length value
$919850631932232$.

The six displayed functions have a nonzero evaluation determinant
$921522$ modulo $1000003$ on the six generic integral symmetric tensors
specified by the seed rule in S-003 below. Hence they are linearly
independent over $\mathbb Q$. As the ambient multiplicity is six, they
form a basis. This finite-field nonzero minor proves characteristic-zero
independence: every denominator in the definitions is invertible at the
chosen prime, so a zero rational minor would reduce to zero there.
No zero modular minor is used to prove an upper rank bound.
\end{proof}

\subsection{S-003. The exact multiplicities $a=6$, $K_3=5$, $q=6$}
For the partition and degree of S-002, extending the partition by two
zero parts to $GL_9$ (plethysm multiplicities are stable once the ambient
vector-space rank contains the partition), one has
\[
 a_\lambda=6,\qquad
 [S_\lambda\mathbb C^9:I(\operatorname{Det}_3)_{12}]=1,
 \qquad K_3(\lambda)=5,\qquad q_\lambda=6.
\]
An explicit highest-weight equation separating the permanent and
 determinant closures is
\[
                  \Psi=2P_0+21P_1+7P_2+35P_3+7P_4.
\]
This verifies independently the two numerical assertions in
Landsberg--Manivel--Ressayre, \emph{Hypersurfaces with degenerate duals
and the Geometric Complexity Theory Program}, Comment.
Math.
Helv.
88 (2013), \S3.2, p.~476, and recomputes both sides of GCT\_KRON F-025.
The original numbers are published results; no priority claim is made.

\begin{proof}[The equation vanishes on the determinant closure]
Write $f(X,Y)=aX^3+3bX^2Y+3cXY^2+dY^3$.
If $H=fQ$ for a binary quartic
$Q=\sum_{j=0}^4\binom4j q_jX^{4-j}Y^j$, its normalized coefficients obey
\[
 h_k=\binom7k^{-1}\sum_{i+j=k}\binom3i\binom4j f_iq_j,
 \qquad(f_0,f_1,f_2,f_3)=(a,b,c,d).
\]
Substituting these expressions in $\Psi$ gives the zero polynomial,
coefficient by coefficient over $\mathbb Q$. This identity is checked
exactly by \texttt{verify\_lmr\_explicit.py}; the identity also follows by
multiplying out the six displayed formulas.

Now let $F(x)=\det_3(\sum_{i=0}^6x_iM_i)$. At a matrix of rank two,
the Hessian of $\det_3$ has rank six. Indeed left and right invertible
changes reduce it to $\operatorname{diag}(1,1,0)$, where the quadratic
term is
\[
              x_{33}(x_{11}+x_{22})-x_{13}x_{31}-x_{23}x_{32},
\]
a quadratic form of rank six. Here the Hessian of $F$ at $Xe_0+Ye_1$ is $6(XS_0+YS_1)$,
so its determinant differs from $H$ only by the nonzero constant $6^7$.
Consequently the Hessian pulled back to
any seven-dimensional linear parameter space has determinant zero at
such a point. For a dense open set of matrix pencils, the binary cubic
$\det_3(XM_0+YM_1)$ is squarefree and its three roots have matrix rank
two. This open set is nonempty: take $M_0=I$ and $M_1$ diagonal with three
distinct entries. Thus $H$ vanishes at every root of $f$ and $f$ divides
$H$. The preceding polynomial identity gives $\Psi(F)=0$ on that dense
open set, hence on every seven-variable determinant projection by
polynomiality.

Every linear endomorphism substitution of $\det_3$ belongs to its orbit
closure, since $GL_9$ is dense in $\operatorname{End}(\mathbb C^9)$.
The same polynomial evaluated on the first seven variables of any
nine-variable translate therefore vanishes; equivalently $\Psi$ lies
in the full determinant orbit-closure ideal. The construction is
highest-weight for $GL_9$ with two terminal zero weights: the last two
simple raising operators act trivially on functions using only the
first seven coefficient coordinates. This proves the required ideal
inclusion.
\end{proof}

\begin{proof}[Exact evaluation certificates for the entire multiplicity space]
All evaluations below are reconstructed from integers by the checker.
For seed $s$, start $u_0=s$, then
\[
 u_{j+1}=1664525u_j+1013904223\pmod{2^{32}},\qquad
 v_j=\left\lfloor u_{j+1}/2^{16}\right\rfloor\pmod{101}.
\]
For determinant/permanent samples, fill the seven $3\times3$ matrices
$M_0,\ldots,M_6$ in row order from the first 63 values. Use
$F_s(x)=\det_3(\sum x_iM_i)$ or
$F_s(x)=\operatorname{per}_3(\sum x_iM_i)$ as appropriate.
The cubic tensor entries are reconstructed by polarizing the six
permutation terms of the $3\times3$ determinant or permanent and dividing
by six. For generic samples, fill the 84 symmetric tensor entries
$T_{ijk}$, $0\leq i\leq j\leq k\leq6$, in lexicographic order from
the same generator.

The prime is $p=1000003$. The seeds, in order, are
$20261122,\ldots,20261127$. Rows are these samples, and columns are
$P_0,P_1,P_2,P_3,P_4,R$. The complete matrices are in
\texttt{results/lmr\_explicit.json}; the checker rebuilds their entries
from the definitions, not from the stored matrices. It gives:
\[
\begin{array}{c|c|c|r}
 \text{family}&\text{rows}&\text{columns}&\text{minor mod }p\\\hline
 \text{generic cubic}&0,1,2,3,4,5&0,1,2,3,4,5&921522\\
 \det_3&0,1,2,3,4&0,1,2,3,5&170120\\
 \operatorname{per}_3&0,1,2,3,4,5&0,1,2,3,4,5&851602
\end{array}
\]
These nonzero minors imply $K_3(\lambda)\geq5$ and $q_\lambda\geq6$
in characteristic zero. The degree-twelve ambient multiplicity is six
by S-002. On the first permanent sample $\Psi$ has value $715852$
modulo $p$, so it is a nonzero ambient highest-weight vector. Its
identical determinant vanishing therefore gives ideal multiplicity at
least one, and $K_3\leq6-1=5$. Combining the bounds proves
$K_3=5$, ideal multiplicity one, and $q=6$. In particular the entire
six-dimensional isotypic multiplicity space has been compared; the
existence of a single equation is not being mistaken for that comparison.
\end{proof}

\paragraph{Grothendieck dictionary.}
Let $X=\mathbb P(\operatorname{Sym}^3(\mathbb C^9)^*)$ and let
$\mathcal I_D$ be the $GL_9$-equivariant coherent ideal sheaf of the
projective determinant orbit closure. The degree-twelve homogeneous
coordinate sequence, restricted to the $S_\lambda$ multiplicity spaces,
is an exact sequence of dimensions $1\to6\to5$. For the permanent
closure the corresponding kernel has dimension zero and the quotient
has dimension six. The quotient here is the image of ambient homogeneous
sections in the coordinate ring; no projective-normality identification
with all sections of $\mathcal O_D(12)$ is required. These are actual ranks of restriction maps of
coherent sections, not merely relations in $K_0$. The explicit six
functions and their evaluation maps are defined over a localization of
$\mathbb Z$ obtained by inverting six and the binomial denominators.
The nonzero minor modulo $1000003$ proves the asserted generic-fiber
ranks after clearing these denominators. It does not assert flatness or
constant ranks in every characteristic. The exterior determinant
constructions belong to the Schur-operation/Cauchy passage of cote~133,
scan~028, and the equivariant module interpretation is the representation
and comodule dictionary of SGA~3, Expos\'e~I, Proposition~4.7.2. No
$\ell$-adic or motivic statement is needed by, or obtained from, these
calculations.

\subsection{S-004. The independent padding variable and the inherited row bounds}
The current prompt uses $z^{n-m}\operatorname{per}_m$ with $z$ independent
of the $m^2$ matrix variables. Several GCT\_KRON proofs explicitly use
internal padding and the row cap $r\leq m^2$. Under the current convention
the elementary support cap is $r\leq m^2+1$. This change does not
invalidate all the inherited bounds, but it must be checked where used.

\paragraph{Statement.}
In the range $n\geq m^{25}$, the following necessary conditions for an
actual obstruction remain valid for independent padding:
\[
 m\geq190,\qquad r\geq52,\qquad
 \lambda_{2j}\geq28-j\ (1\leq j\leq26),\qquad
 \lambda_3\geq557\ \text{or}\ \lambda_{52}\geq3.
\]
The two-order-ideal construction of M-019, with rectangular root sets,
also proves
\[
              \lambda_4\geq27\quad\text{or}\quad
              \lambda_2\geq238356708.
\]
For $r\leq m^2$ the inherited threshold $239063403$ still applies.
For the extra possible row $r=m^2+1$, the inherited proof does not justify
that stronger numerical threshold. This is a correction of the proof's
scope, not a counterexample to the stronger inequality.

\begin{proof}
F-009 already gives a homogeneous determinantal representation with an
independent auxiliary variable $L$; hence its $m\geq190$ consequence
is unchanged. M-008 represents every $L^{n-m}f$ in a fixed number of
variables and does not use the permanent's row cap. Its $r\geq52$
consequence is likewise unchanged.

For the c=25 staircase in M-014, all cuts other than $k=2$ have $k\geq4$.
Their suffix variable count is $R=r-k+1\leq m^2-2$. The same estimates
in that proof therefore apply with its numerical auxiliary bound $m^2$.
At $k=2$ use M-013 with $\ell=12$, $b=26$ and
$h=\lceil\sqrt{r-1}\rceil\leq m$. Its required size is at most
$2m\binom{m^2+12}{12}$. The ratio of this to $m^{25}$ is
\[
               \frac2{12!}\prod_{i=1}^{12}(1+i/m^2),
\]
which decreases with $m$ and is strictly less than one at $m=190$ by
exact integer comparison. Thus $\lambda_2\leq26$ is excluded and the
same staircase follows.

For A-025 replace the last-group cap $m^2-51$ by $m^2-50$ in its
pruned cover. Explicitly, put
\[
 A(m,v)=(m+1)\binom{49+t}{49}-49\binom{49+t}{50},\qquad t=\min(m,v).
\]
With $u=360,v=195$, the modified count is
\[
 T^{\rm ext}_*(m)=\binom{49+\min(360,m)}{49}
       +A(m,195)+(m^2-50)A(m-1,195).
\]
The standalone checker verifies
$25T^{\rm ext}_*(m)<24m^{25}$ for all 171 integers $190\leq m\leq360$.
For $m\geq360$, write $C=\binom{409}{49}$,
$D=\binom{244}{49}$ and $c=1911/10$. The normalized count is the sum of
\[
 C/m^{25},\qquad D(m+1-c)/m^{25},\qquad
                 D(m^2-50)(m-c)/m^{25}.
\]
Each term is positive and decreasing. For the last term its logarithmic
derivative is at most
$3/m+3/m-25/m<0$; the same inequalities used by A-025 remain valid,
with 50 in place of 51. Thus the checked endpoint establishes the
infinite tail. The same support cover and S-001 then give the same
$557$-or-$3$ necessary condition.

Finally, in M-019 use $m$ inward root trees and $m+1$ outward root trees.
They offer $m(m+1)\geq m^2+1\geq r$ distinct bridges, enough for all
$r-1$ variable labels and the constant label. Every path and cofactor
argument in M-019 is unchanged. With $\ell=12$ each tree has at most
$Q=\binom{m^2+12}{12}$ vertices. If the attached axis ideals both have
$t+1$ members, the required size is at most
\[
                         (2m+1)(t+1)Q.
\]
The function
\[
 \frac{m^{25}}{(2m+1)\binom{m^2+12}{12}}
       =\frac{12!}{(2+1/m)\prod_{i=1}^{12}(1+i/m^2)}
\]
is increasing on $m>0$. At $m=190$ its integer floor is exactly
$J=238356682$, as checked by the two integer inequalities defining a
floor. Set $t=J-1$. If $\lambda_4\leq26$ and
$\lambda_2\leq26+t=238356707$, the support is covered by these two
axis ideals plus the order-26 simplex, exactly as in M-019. The required
matrix size is at most $m^{25}\leq n$, so the whole ambient highest-weight
space survives on the determinant. This forbids an obstruction and
proves the asserted alternative.
\end{proof}

\paragraph{Other consequences of the convention.}
A-035's basic partial-Hessian obstruction exclusion is unaffected:
$m^2+1<n^2$ for every $n>m\geq2$, so the essential dimension remains
proper. Its uniform guarantee that \emph{every} relevant partial-Hessian
binary coefficient algebra is covered by the $r\leq2n$ injection changes
from $m^2\leq2n$ to $m^2+1\leq2n$. The determinant-only support and jet
criteria F-016, M-013, M-015--M-019 are stated in terms of the actual row
count $r$ and retain those exact hypotheses. Internally padded permanent
families used for lower bounds remain available inside the independently
padded closure by specializing the separate padding variable to a matrix
entry, so such lower bounds are not lost. A claim asserting equality or
an upper bound from the smaller essential dimension must be checked
separately, as above.

\subsection{The F-016 dependency inventory}
The complete inventory with exact per-statement proof locators and
subclaim scope is in \texttt{dependencies.tsv}; excluded cross-references
are in \texttt{nondependencies.tsv}. Direct support applications occur
in F-016's own all-copy consequence, F-017, F-018, F-021--F-023,
F-025--F-026, M-011, M-013, M-015--M-019, A-025, A-031,
O-015.b, O-018.a, O-020, O-021 and O-022. The following additional
statements inherit those applications: M-014, O-015.c, O-018.b,
O-023--O-025 and O-027; O-027 also explicitly repeats the support step.
For O-024 the support interpretation is separate from its combinatorial
cover identity. O-026's multiplicity interpretation depends on the
support lemma, whereas its exact combinatorial minima do not.
A-035 uses it only in the optional O-027 full-ambient-coverage corollary,
not in the basic bracket nonvanishing or binary-algebra injection.
O-031, QA033 and QA033-CB use it only for their conditional eventual
multiplicity consequence; their unresolved Jacobian gate is not supplied
by the lemma. M-020 originally inherits it for the illustrative F-025
example, not for its general saturation statement. The specialized
F-025 example now has the independent proof S-003, so that numerical
comparison no longer needs F-016.

Each inventory check confirms the particular inference from coefficient
support to an all-copy ideal or quotient statement. It does not constitute
an independent proof of every other geometric or numerical input to the
referenced finding. In every direct occurrence the same valid step is
used: a highest-weight polynomial belongs to the polynomial algebra on
the indicated suffix-bounded coefficient set. Dominant projection to
that set makes its pullback injective; alternatively an exact determinant
representation of the retained padded forms makes every determinant-ideal
highest-weight vector vanish on the permanent family. Both arguments
apply to all vectors of the multiplicity space. Numerical substitutions
using the row cap must use the corrected convention in S-004.

\subsection{S-005. A finite 804-dimensional gate for the $(8,9)$ question}
This is a proved reduction of one concrete four-row family, not an
explicit obstruction. Let $F$ be a degree-$n$ form in four variables and
write $F(v)=T(v,\ldots,v)$ with symmetric $n$-linear tensor $T$.
Distinguish $e_0$ and put $U=\langle e_1,e_2,e_3\rangle$. For $v\in U$
write
\[
 F_j(v)=T(e_0^{n-j},v^j),\qquad 0\leq j\leq n.
\]
Thus $F_0$ is a scalar. Let $\mathcal H_d(n,c)$ be the ambient
highest-weight space of weight $(nd-3c,c,c,c)$ in outer degree $d$.
Assume $1\leq c\leq n$. Define
\[
 \mathcal R_c=
 \left[\operatorname{Sym}\left(\bigoplus_{j=2}^c
                        \operatorname{Sym}^j\mathbb C^3\right)\right]_{\det^c},
\]
where the bracket denotes the $GL_3$ determinant-character space, so
all its elements have total ordinary $GL_3$ weight $3c$.

\paragraph{Uniform stabilization.}
For every $d\geq3c$, multiplication by $F_0^{d-3c}$ is an isomorphism
\[
             \mathcal H_{3c}(n,c)\longrightarrow\mathcal H_d(n,c),
             \qquad\mathcal H_{3c}(n,c)\simeq\mathcal R_c.
\]
The assertion concerns the actual polynomial highest-weight spaces,
not only equality of their dimensions.

\begin{proof}
Localize at $F_0\ne0$ and define the centered coefficients
\[
 C_j(v)=\sum_{i=0}^j(-1)^{j-i}\binom ji
               F_i(v)\frac{F_1(v)^{j-i}}{F_0^{j-i}}
       =T\left(e_0^{n-j},
          \left(v-\frac{F_1(v)}{F_0}e_0\right)^j\right).
\]
Then $C_1=0$, and $F_0^{j-1}C_j$ is polynomial of coefficient degree $j$.
For the terms $i=0,1$, the apparent denominator $F_0^j$ is reduced by
the factor $F_0$ in $F_i$ or by collecting the two terms; for $i\geq1$
the displayed denominator has degree at most $j-1$ directly.
More simply, the $i=0$ term is $(-1)^jF_1^j/F_0^{j-1}$, and all other
terms have denominator at most $F_0^{j-1}$.

The three-dimensional unipotent group $g(e_i)=e_i+a_ie_0$, $g(e_0)=e_0$,
translates $F_1/F_0$ freely and fixes every $C_j$. The change between
$(F_0,F_1,F_2,\ldots,F_c)$ and
$(F_0,F_1,C_2,\ldots,C_c)$ is triangular and invertible after localizing
at $F_0$. Consequently its invariant ring is exactly
\[
             \mathbb C[F_0,F_0^{-1},C_2,\ldots,C_c].
\]
By S-001, every vector of $\mathcal H_d(n,c)$ uses only the coefficients
$F_0,\ldots,F_c$, since every coefficient generator has tail degree at
most $\lambda_2=c$. The remaining highest-weight condition with tail
$(c,c,c)$ says precisely that its centered expression is $SL_3$ invariant
with $GL_3$ character $\det^c$.

Decompose such an invariant $I$ into pieces $I_k$ with $k$ centered-jet
factors. The weighted degree condition is $\sum j=3c$ in each monomial,
and $j\geq2$, so $k\leq3c/2$. Its denominator is at worst
$F_0^{3c-k}$. Hence
\[
                  L(I)=\sum_kF_0^{3c-k}I_k(C_2,\ldots,C_c)
\]
is a polynomial of coefficient degree $3c$. Under
$\operatorname{diag}(t_0,g)$, $C_j$ transforms as
$t_0^{n-j}\operatorname{Sym}^j g$, so $L(I)$ has character
$t_0^{3cn-3c}\det(g)^c$ and is highest-weight. At $F_0=1,F_1=0$ it
recovers $I$ on freely assignable $C_2,\ldots,C_c$, proving injectivity.

Conversely the localized expression of a degree-$d$ highest-weight
polynomial is $\sum_kF_0^{d-k}I_k(C)$. This follows from coefficient
homogeneity and the tail character, and it has no additional dependence
on the freely translated $F_1/F_0$. It equals $F_0^{d-3c}L(I)$.
For $d\geq3c$ this is a polynomial identity; no division by $F_0$ is
being asserted in the unlocalized coordinate ring. The same argument
shows that multiplication by $F_0^{3c-d}$ embeds every earlier
$\mathcal H_d(n,c)$ into $\mathcal H_{3c}(n,c)$. This proves all claims.
\end{proof}

\paragraph{The concrete dimension.}
At $n=9,c=8$ one obtains
\[
          \dim\mathcal H_{24}(9,8)
                  =a_{(192,8,8,8)}=804.
\]
The standard-library checker \texttt{check\_centered\_jet\_gate.py}
performs two independent character computations. The centered-jet
computation takes the coefficient of the determinant character using
six $GL_3$ Weyl terms in
\[
          \prod_{\substack{\alpha\in\mathbb Z_{\geq0}^3\\
                           2\leq|\alpha|\leq8}}(1-x^\alpha)^{-1}.
\]
Its signed dimensions are
\[
       2604128-2468696-2468696+2199430+2237895-2103257=804.
\]
The second computation applies all 24 $GL_4$ Weyl terms directly to
$\operatorname{Sym}^{24}(\operatorname{Sym}^9\mathbb C^4)$ at
$(192,8,8,8)$ and gives the same value. In these terms the total tail
weight is at most 24; each monomial in nonconstant tail generators
therefore uses at most 24 factors and has a unique completion to outer
degree 24 by powers of the pure first-coordinate generator. This
justifies the checker's three-dimensional integer knapsack for the
direct plethysm calculation.

\paragraph{The padded side is completely known in this family.}
For $(m,n)=(8,9)$, CUE\_GCT R-13/R-14 identify the four-variable padded
image with all forms $\ell q$, where $q$ is an arbitrary octic.
It contains the complete linear subspace $e_0^*\operatorname{Sym}^8V^*$.
Its coefficients $F_0,\ldots,F_8$ vary freely, and S-001 puts every
vector of $\mathcal H_d(9,8)$ in their polynomial algebra. Restriction
to this subspace is therefore injective on the entire highest-weight
space, proving
\[
                  q_{(9d-24,8,8,8)}=a_{(9d-24,8,8,8)}
\]
whenever the displayed tuple is a partition. In particular $q=804$
at $d=24$.

\paragraph{One finite rank decides all degrees of this family.}
The determinant coordinate ring is a domain and $F_0$ is nonzero on it.
Thus an ideal vector in degree $d>24$ factors as
$F_0^{d-24}$ times one in degree 24 by the proved polynomial identity,
and that factor is again in the determinant ideal. A nonzero ideal
vector in an earlier degree can be multiplied by $F_0^{24-d}$.
Therefore a gap exists in this family in \emph{some} degree if and only
if the restriction map from the explicit 804-dimensional space
$\mathcal H_{24}(9,8)$ to the determinant coordinate ring has rank below
804. This reduction resolves the all-degree issue without asserting any
uncomputed determinant-side rank.

\paragraph{A rank-803-or-804 dichotomy, using this run's jet certificate.}
Let $Y$ be the closure of the determinant image on the 165 coefficients
$F_0,\ldots,F_8$. The root certificate
\texttt{root/results/det9\_nested\_profiles.json}, profile $(8,8,7)$,
has rank 164 and minor 76 modulo 101. Combined with the determinant
dimension upper bound 164, it makes $Y$ an irreducible hypersurface.
The same full-rank subprojection contains every coefficient through
order seven, proving dominance of that smaller jet.

Since the line-stabilizer parabolic preserves the 165-coordinate
algebra, it preserves the principal prime ideal of $Y$. Its irreducible
homogeneous generator $G$ is therefore a parabolic semi-invariant of
weight $(9D-3c_G,c_G,c_G,c_G)$ for some nonnegative integer $c_G$.
If $c_G\leq7$, S-001 places $G$ in the dominant order-seven jet algebra,
a contradiction. Hence $c_G\geq8$.

If a nonzero vector of $\mathcal H_{24}(9,8)$ is divisible by $G$, its
polynomial quotient is a highest-weight semi-invariant of tail weight
$(8-c_G,8-c_G,8-c_G)$. Negative tail weights are impossible in the
polynomial algebra, so $c_G>8$ gives no kernel. If $c_G=8$, the quotient
has tail weight zero and is a power of $F_0$ times a scalar. Moreover
$G$ has coefficient degree $D\leq24$: otherwise the stabilization
identity would make $G$ divisible by $F_0$, contradicting irreducibility
and the fact that $Y$ is not the hyperplane $F_0=0$. Thus the kernel
has dimension exactly one when $c_G=8$, and zero when $c_G>8$.
Consequently the two possible outcomes of the finite calculation are
\[
  q=804>K_9=803\quad(c_G=8),\qquad
  q=K_9=804\quad(c_G>8).
\]
The value of $c_G$, the 804-dimensional determinant evaluation rank,
and the defining polynomial $G$ have not been computed here. This
corollary is a strictly smaller gate for this chosen direction, not a
claim that either alternative has been decided.

\paragraph{Geometric dictionary.}
Centering is an explicit coordinate calculation on the principal affine
open $D(F_0)$ and a quotient by a split unipotent group. Clearing powers
of $F_0$ constructs actual homogeneous sections. The final rank test is
the $S_{(192,8,8,8)}$ part of the equivariant coherent ideal of the
quaternary determinant image; the principal ideal $(G)$ describes its
projection to the octic-polar coefficient scheme. The all-degree
reduction uses localization and the fact that multiplication by a
nonzero section on an integral scheme is injective. It does not replace
the coordinate ring by its normalization or infer ranks from its
$K_0$ class. The representation/comodule interpretation is again SGA~3,
Expos\'e~I, Proposition~4.7.2, with Schur operations as in the supplied
cote~133, scan~028. No assertion about weights of $\ell$-adic cohomology
or motives is needed by this affine calculation.
